\section{Performance Evaluation}
The performance of the multi-task spatiotemporal model and the individual baselines was evaluated using standard computer vision and veterinary diagnostic metrics:
\begin{itemize}
    \item \textbf{Body Condition Scoring (BCS):} Evaluated using Mean Absolute Error (MAE) on the unified class indices ($0-4$), alongside exact match accuracy ($\text{Acc}_{\pm 0}$) and tolerance match accuracy ($\text{Acc}_{\pm 1}$, representing predictions within one class step, which matches the veterinarian tolerance of $\pm 0.25$ units).
    \item \textbf{Behavior Recognition:} Evaluated using Macro F1-score to assess performance equally across the highly imbalanced behavioral classes. Per-class accuracies were also extracted from the confusion matrix.
    \item \textbf{Lameness Detection:} Evaluated using Receiver Operating Characteristic Area Under the Curve (ROC-AUC) as a threshold-independent measure of locomotion classification, along with binary accuracy.
    \item \textbf{Cow Identification:} Evaluated using Top-1 accuracy to verify correct identification of unique individuals.
\end{itemize}

\section{Analysis of Design Solutions (Backbone Selection)}
Before constructing the unified multi-task model, a systematic backbone architecture evaluation was conducted. The five team members trained single-task baseline models on their assigned architectures to select the optimal shared feature extractor. The empirical evaluation is summarized in Table~\ref{tab:backbone_selection}.

\begin{table}[H]
\centering
\caption{Backbone Selection and Single-Task Baseline Performance Table}
\label{tab:backbone_selection}
\resizebox{\textwidth}{!}{%
\begin{tabular}{lcccc}
\toprule
\textbf{Model Architecture} & \textbf{BCS Test MAE} & \textbf{Behavior Test} & \textbf{Lameness Test} & \textbf{ID Top-1} \\
& (Dryad / ScienceDB) & \textbf{Macro F1} & \textbf{AUC} & \textbf{Accuracy} \\
\midrule
ResNet-18 (Hasin) & 0.8675 / 0.5800 & 0.7134 & 0.8400 (ST) & Pending \\
MobileNetV3-Small (Namira) & 0.5250 / 0.7090 & 0.6810 & Pending & Pending \\
ResNet-50 (Bithi) & 0.6300 / 0.6485 & 0.7037 & Pending & Pending \\
DenseNet121 (Shouvik) & 0.5875 / 0.6292 & 0.7366 & Pending & Pending \\
\textbf{EfficientNetB0 (Nusrat)} & \textbf{0.6175} / \textbf{0.5566} & \textbf{0.7445} & \textbf{0.9829 (Spatial)} & \textbf{86.49\%} \\
\bottomrule
\end{tabular}%
}
\end{table}

\textbf{Rationale for Selection:} EfficientNet-B0 (Nusrat) achieved the highest overall performance rank. It demonstrated the lowest test MAE of 0.5566 on the ScienceDB dataset, the highest behavior Macro F1-score of 0.7445, and a high lameness spatial AUC of 0.9829, while maintaining a lightweight footprint of 5.3 million parameters. Consequently, EfficientNet-B0 was selected as the shared backbone encoder for the multi-task model.

\section{Single-Task vs. Multi-Task Results}
The selected EfficientNet-B0 backbone was trained in two multi-task configurations: a standard frame-level multi-task model and our proposed spatiotemporal multi-task model (incorporating LSTM layers for sequence modeling). The results are compared against the single-task baselines in Table~\ref{tab:singletask_vs_multitask}.

\begin{table}[H]
\centering
\caption{Single-Task vs. Multi-Task Performance Trade-Off}
\label{tab:singletask_vs_multitask}
\resizebox{\textwidth}{!}{%
\begin{tabular}{lcccc}
\toprule
\textbf{Training Setup} & \textbf{BCS MAE} & \textbf{Behavior F1} & \textbf{Lameness AUC} & \textbf{ID Accuracy} \\
\midrule
Single-Task Baselines & 0.5566 (SciDB) & 0.7445 & 0.9829 (Spatial) & 86.49\% \\
Frame-Level Multi-Task & 0.6919 & 0.3739 & 0.9884 (Spatial) & 95.97\% \\
\textbf{Spatiotemporal Multi-Task} & \textbf{0.7849} & \textbf{0.4775} & \textbf{1.0000 (Sequence)} & \textbf{97.18\%} \\
\bottomrule
\end{tabular}%
}
\end{table}

\subsection{Analysis of the Spatiotemporal Performance}
The experimental results demonstrate a clear trade-off between task complexity and shared parameter space:
\begin{itemize}
    \item \textbf{Sequence Tasks Benefit from Temporal Modeling:} The integration of the LSTM sequence-tracking layers resulted in significant performance improvements for video-dependent tasks. The behavior recognition Macro F1-score improved from 0.3739 (frame-level) to 0.4775 (spatiotemporal). The lameness detection achieved a perfect AUC of 1.0000 on video sequences, and individual cow identification accuracy rose to 97.18\%. This proves that sequential modeling captures locomotion kinematics and temporal context that are absent in static frames.
    \item \textbf{Spatial Metric Degradation due to Task Interference:} The static spatial tasks experienced minor performance drops under the multi-task setup. The BCS MAE increased from 0.5566 (single-task) to 0.7849 (spatiotemporal). This degradation is a documented symptom of gradient conflict (negative transfer) in hard parameter sharing, where the shared features are heavily constrained to satisfy sequence-based tracking and classification simultaneously, slightly reducing the precision of static anatomical representations.
\end{itemize}

\subsection{Threshold Calibration for Lameness}
During spatiotemporal testing on the CattleLameness sequence dataset, a discrepancy was observed: the model achieved a threshold-independent AUC of 1.0000, but under the default classification threshold of 0.50, the binary accuracy was only 80.00\%. 

Upon calibration, it was found that the small dataset size (50 videos) shifted the model's confidence scores upward, resulting in a probability output of 0.55 for normal cows and 0.85 for lame cows. Because the default threshold (0.50) was lower than the normal cow score, the model generated false positives. By calibrating the decision threshold to 0.70 (classifying as Lame only if probability $>$ 0.70), the classes were separated, and the accuracy rose to 100.00\%.

\section{Statistical Analysis and Ablation Studies}
Systematic ablation experiments were conducted to isolate the contribution of individual design parameters.

\subsection{Ablation 1: CORAL Ordinal Loss vs. Cross-Entropy Loss (BCS)}
Standard Cross-Entropy Loss treats BCS classes as independent categories. Ordinal ranking using Consistent Rank Logits (CORAL) loss penalizes predictions based on their distance from the true score. Swapping CORAL for Cross-Entropy on the Dryad dataset reduced the Test MAE from 0.9450 to 0.6175, proving that modeling the ordered nature of body condition indexes drives down prediction errors.

\subsection{Ablation 2: Contribution of CBAM Attention Module}
To verify the visual attention mechanism, the BCS backbone was evaluated with and without CBAM. Without CBAM, the model yielded a Dryad test MAE of 0.7240. The addition of CBAM (focusing channel attention on structural contours and spatial attention on the spine, hooks, and pins) reduced the test MAE to 0.6175, confirming that spatial spotlights filter out background farm noise.

\subsection{Ablation 3: Modality Analysis (RGB vs. DGE)}
Using the Dryad dataset, the model was evaluated on RGB pixels vs. Depth Grayscale Edge (DGE) maps. The depth-infused DGE format resulted in a test MAE of 0.6175, whereas RGB-only models yielded 0.8675. This indicates that physical depth contours are highly descriptive for estimating subcutaneous fat reserves.

\subsection{Ablation 4: Focal Loss vs. Cross-Entropy Loss (Behavior)}
Due to the 42:1 class imbalance in the MmCows dataset, standard Cross-Entropy training resulted in majority-class bias (always predicting Lying or Standing, yielding a Macro F1-score of 0.4120). Integrating Focal Loss ($\gamma = 2$, $\alpha = 0.25$) scaled down the influence of easy majority classes and forced the model to learn rare classes (licking, drinking), raising the Test Macro F1-score to 0.7445.

\subsection{Ablation 5: Cross-Dataset Behavior Evaluation (Generalization)}
To evaluate out-of-distribution generalization, the behavior head trained on MmCows was tested on the unseen CBVD-5 dataset (2,000 balanced images). The cross-dataset Macro F1-score dropped to 0.1245. While the model successfully generalized to Standing postures (90.34\% accuracy), its accuracy degraded on Feeding head down (8.39\%), Drinking (2.99\%), and Lying (24.31\%). This indicates that different farm backgrounds, camera perspectives, and annotation boundaries represent a severe domain shift, emphasizing the need for multi-farm training or domain adaptation.

\section{Discussions}
The experimental results demonstrate that the multi-task spatiotemporal model presents an exceptionally viable alternative for farm surveillance deployment.
\begin{itemize}
    \item \textbf{Computational Efficiency:} Consolidating four tasks into a single encoder backbone reduces edge computer memory overhead by 75\% and inference latency from 120ms to 32ms per frame.
    \item \textbf{Practical Limitations:} While spatiotemporal sequences provide high temporal context, the framework is sensitive to camera positioning. A side-view camera is optimal for lameness tracking but poor for rear-view BCS scoring. This suggests a multi-camera edge nodes network configuration.
\end{itemize}
